% Compile with LuaLaTeX: lualatex when_every_parenthesis_matters_AMM_note_v3_ja_lualatex.tex
\documentclass[11pt]{ltjsarticle}

\usepackage[margin=0.92in]{geometry}
\usepackage{amsmath,amssymb,amsthm}
\usepackage{booktabs}
\usepackage[hidelinks]{hyperref}

\newtheorem{theorem}{定理}
\newtheorem{lemma}[theorem]{補題}
\newtheorem{proposition}[theorem]{命題}
\newtheorem{corollary}[theorem]{系}

\renewcommand{\proofname}{証明}
\renewcommand{\abstractname}{概要}
\renewcommand{\refname}{参考文献}
\renewcommand{\tablename}{表}

\newcommand{\Q}{\mathbb{Q}}
\newcommand{\Rcal}{\mathcal{R}}
\newcommand{\Zcal}{\mathcal{Z}}
\newcommand{\Pcal}{\mathcal{P}}
\newcommand{\Bcal}{\mathcal{B}}

\title{すべての括弧が必要になるとき}
\author{栗山貴之}
\date{}

\begin{document}
\maketitle

\begin{abstract}
変数の順序を $x_1,\ldots,x_n$ に固定し、$+$、$\times$、$/$ を用いて完全括弧付き式を作る。
最外側以外の任意の一組の括弧を取り除くと、その結果得られる中置式が表す有理関数が変化するとき、
その式を\emph{剛直}であるという。三行からなる局所規則によって、剛直な式の完全な構造分類が得られる。
さらに、相異なる剛直な木は、実際に相異なる有理関数を表す。剛直な式の個数を $r_n$ とすると、
\[
 r_2=3,
 \qquad
 r_n=2F_{n+1}\quad(n\ge3),
\]
となる。ここで $F_1=F_2=1$ はフィボナッチ数である。
\end{abstract}

\section{問題}

括弧は、見た目には存在していても、数学的には何の役割も果たしていないことがある。たとえば、
\[
 x_1+(x_2/x_3)=x_1+x_2/x_3
\]
である。一方、
\[
 x_1/(x_2+x_3)\ne x_1/x_2+x_3
\]
である。では、内側にあるどの括弧も無意味ではない完全括弧付き式は、いくつあるだろうか。

この順序で並ぶ相異なる不定元 $x_1,\ldots,x_n$ を固定する。
\emph{完全括弧付き式}とは、葉が左から右へ $x_1,\ldots,x_n$ と並び、内部頂点に
\[
 +,\qquad \times,\qquad /
\]
のいずれかが付された平面完全二分木である。式は $\Q(x_1,\ldots,x_n)$ の元として比較する。
式全体を囲む一組の括弧は考慮しない。なぜなら、それを取り除いても構文解析も表される有理関数も変化せず、
これまで数えることにすると剛直な式が一つも存在しなくなるからである。
残る $n-2$ 組の括弧は、それぞれ二つの内部頂点を結ぶ一辺に対応する。

一組の括弧を取り除いた後に得られる部分括弧付き中置式は、次の規約に従って構文解析する。
$\times$ と $/$ は同じ優先順位をもち、左から右へ評価する。一方、$+$ はそれらより低い優先順位をもち、
これも左結合とする。答えはこの規約に依存する。たとえば、積と商を右結合として構文解析するなら、
$(A/B)\times C$ における括弧は一般に必要になる。

他のすべての括弧を残したまま、その一組の括弧だけを取り除き、上の規約に従って結果の式を構文解析したとき、
式が表す有理関数が変化するなら、その括弧を\emph{必要}であるという。
$n-2$ 個の候補となる括弧がすべて必要であるとき、その式を\emph{剛直}であるという。
葉を少なくとも二つもつ式を\emph{非自明}であるという。
$x_1,\ldots,x_n$ 上の剛直な式の個数を $r_n$ とする。

$n=2$ のとき、候補となる括弧は一つもない。したがって、一つの演算だけからなる三つの式は、
すべて空虚な意味で剛直である。よって $r_2=3$ であり、これが定理~\ref{thm:main} における
例外的な初期値の理由である。変数が三つの場合、剛直な式は次の六つである。
\[
\begin{gathered}
 (x_1+x_2)\times x_3,\qquad
 (x_1+x_2)/x_3,\qquad
 x_1\times(x_2+x_3),\\
 x_1/(x_2+x_3),\qquad
 x_1/(x_2\times x_3),\qquad
 x_1/(x_2/x_3).
\end{gathered}
\]
本稿全体を通じて、最外側の一組の括弧は表示していない。より大きな例として、
\[
 ((x_1+x_2)/(x_3/(x_4+x_5)))
\]
は剛直である。最外側の括弧を除けば、残る三組の括弧はすべて必要である。

本稿の対象は、括弧が最小限に付された中置文字列とは区別されるべきである。
本稿では完全括弧付き二分木から出発し、あらかじめ指定された各内側の括弧が、一つずつ必要であるかを問う。
基礎となる木を固定することによって、括弧の必要性が辺に局所的な性質になる。
最初は演算ラベルを付した木を数えるが、命題~\ref{prop:distinct} により、二つの剛直な木が
同じ有理関数を表すことはないと分かる。

本稿では三つの演算だけを扱う。減算を加えると、別の数え上げ問題が生じる。

\section{局所規則}

根の演算が $+$ である内部子を\emph{加法的}であるという。
次の表は、ある子を囲む括弧が必要であるかどうかを決定する。

\begin{table}[ht]
\centering
\small
\begin{tabular}{ccc}
\toprule
親 & 左の内部子 & 右の内部子 \\
\midrule
$+$      & 決して必要でない & 決して必要でない \\
$\times$ & 加法的な場合に限り必要 & 加法的な場合に限り必要 \\
$/$      & 加法的な場合に限り必要 & 常に必要 \\
\bottomrule
\end{tabular}
\caption{内部子を囲む括弧が必要となる条件。}
\label{tab:local}
\end{table}

\begin{lemma}[局所規則]\label{lem:local}
完全括弧付き式の根でない内部頂点を $v$ とし、その親を $p$ とする。
辺 $pv$ に対応する一組の括弧が必要であるための必要十分条件は、
表~\ref{tab:local} の対応する欄で「必要」とされていることである。
\end{lemma}

\begin{proof}
考察している一組以外の括弧はすべて残す。局所的には、式は三つの空でないブロック $A,B,C$ からなる。

まず、任意の空でないブロックを、指定された任意の正の有理数に特殊化でき、しかもそのすべての部分式の値を
正にできることに注意する。これは再帰的に分かる。加法節点では二つの子の値を $(q/2,q/2)$ とし、
乗法節点または除法節点では $(q,1)$ とすればよい。
$A,B,C$ は互いに素な変数集合を含むので、それぞれの値を独立に指定できる。

親が $+$ であるとする。どちらの子を露出させても値は変化しない。
加法的な子は結合則によって吸収され、乗法または除法を根にもつ子は優先順位によって保護されるからである。

次に、親が $\times$ であるとする。根が $\times$ または $/$ である子は、括弧を取り除いて露出させてもよい。
実際、
\[
 (A\times B)\times C=A\times B\times C,
 \qquad
 (A/B)\times C=A/B\times C
\]
であり、右の子についても同様の恒等式が成り立つ。一方、加法的な子は露出させることができない。たとえば、
\[
 (1+1)\times2=4\ne3=1+1\times2
\]
である。加法的な右の子についても、同様の正の特殊化を用いればよい。

最後に、親が $/$ であるとする。左の子の根が $\times$ または $/$ なら、左結合の規約により括弧を取り除ける。
一方、左の子が加法的なら括弧を取り除くことはできない。たとえば、
\[
 (1+1)/2=1\ne\frac32=1+1/2
\]
である。内部頂点である右の子は、根の演算が何であっても括弧を必要とする。
三つの可能な根演算は、たとえば次のように区別される。
\[
 \frac{1}{1+1}\ne \frac11+1,
 \qquad
 \frac{1}{1\times2}\ne \frac11\times2,
 \qquad
 \frac{1}{1/2}\ne \frac11/2.
\]
先ほどの正の特殊化に関する観察によって、これらの数値的な証人を任意の大きさのブロックの内部で実現できる。
したがって、表~\ref{tab:local} で必要と判定されるすべての場合において、括弧を取り除くと、
影響を受ける局所部分木が表す有理関数が変化する。

残るのは、この局所的な変化が木の上方で消えないことを示すことである。
$K=\Q(x_1,\ldots,x_n)$ とし、括弧を取り除く前後に、影響を受ける部分木が表す二つの有理関数を
$f,g\in K$ とする。先ほどの特殊化から $f\ne g$ である。
両側に現れるすべての中間式、およびすべての兄弟部分木は、$K$ の非零元である。
実際、すべての変数を $1$ とおけば正の値をもつ。
各祖先において、兄弟部分木が表す有理関数 $h$ を固定すると、現在の値 $u$ は次のいずれかで変換される。
\[
 u\longmapsto u+h,
 \qquad
 u\longmapsto hu,
 \qquad
 u\longmapsto u/h,
 \qquad
 u\longmapsto h/u.
\]
ここで $h\ne0$ であり、最後の場合には $u\ne0$ である。
これらの写像は、それぞれの定義域上で単射である。
したがって、根に至る道に沿った帰納法により、二つの完全な式は相異なる有理関数を表す。
\end{proof}

\begin{corollary}[文脈独立性]\label{cor:context}
$T$ を完全括弧付き式 $E$ の部分木とし、$\pi$ を $T$ の内部にある候補括弧とする。
このとき、$\pi$ が $T$ において必要であるための必要十分条件は、$E$ において必要であることである。
\end{corollary}

\begin{proof}
$\pi$ を取り除いても $T$ が表す有理関数が変化しないなら、式全体も変化しない。
逆に、$\pi$ を取り除くことで $T$ が表す関数が変化するなら、補題~\ref{lem:local} の証明における
単射性の議論により、その変化はすべての祖先を通って根まで伝わる。
\end{proof}

\section{構造とフィボナッチ数による数え上げ}

局所規則から、完全な構造記述が得られる。

\begin{proposition}[構造分類]\label{prop:structure}
少なくとも二つの葉をもつ完全括弧付き式が剛直であるための必要十分条件は、
次の三つのうちちょうど一つが成り立つことである。
\begin{enumerate}
\item 根が $+$ であり、二つの子がともに葉である。
\item 根が $\times$ であり、各子が、葉または二葉の和のいずれかである。
\item 根が $/$ であり、左の子が葉または二葉の和のいずれかで、右の子が葉または剛直な式のいずれかである。
\end{enumerate}
\end{proposition}

\begin{proof}
まず、式が剛直であるとする。根が $+$ なら、局所規則によって内部子は一つも許されないので、
二つの子はともに葉である。根が $\times$ なら、すべての内部子は加法的でなければならない。
系~\ref{cor:context} により、そのような子自身も剛直であり、第一の場合から、
二つの葉を $+$ で結んだ式でなければならない。根が $/$ なら、左の子については同じ議論が成り立つ。
一方、局所規則は内部右子の根演算には制限を課さない。その右部分木は、系~\ref{cor:context} により剛直である。

逆に、列挙した三つの形のいずれかをもつ式を考える。根の直下にある各括弧は局所規則により必要である。
また、剛直な子の内部にある各括弧は、系~\ref{cor:context} によって式全体の中でも必要である。
したがって、式全体が剛直である。
\end{proof}

この構造分類から、木の個数を数えたときに $\Q(x_1,\ldots,x_n)$ の中で隠れた衝突が生じないことも分かる。

\begin{proposition}\label{prop:distinct}
相異なる剛直な式は、相異なる有理関数を表す。
\end{proposition}

\begin{proof}
$x_i$ および $x_i+x_{i+1}$ を\emph{基本ブロック}と呼ぶ。
命題~\ref{prop:structure} により、任意の葉または剛直な式が表す有理関数は、
\[
 \prod_{j=1}^k b_j^{\varepsilon_j},
 \qquad \varepsilon_j\in\{1,-1\}
\]
と書ける。ここで $b_j$ は、互いに素な連続区間上に台をもち、左から右へ並べると式の変数全体を分割する
基本ブロックである。実際、乗法は二つの正の基本ブロックを連結し、除法は左の基本ブロックを最初に置き、
右の子の因数分解に現れるすべての符号を反転させる。

各基本ブロックは、一意分解整域 $\Q[x_1,\ldots,x_n]$ におけるモニックな既約一次多項式であり、
相異なる基本ブロックは互いに同伴ではない。したがって、有理関数から、符号付き基本因子の列が一意に定まる。

あとは、その列から木を復元すればよい。一因子だけからなる列は、葉 $x_i$ または二葉の和 $x_i+x_{i+1}$ を与える。
二つの正の因子からなる列は、その二つの基本ブロックの積を与える。
それ以外のすべての場合、構造分類により、第一因子は正で第二因子は負でなければならない。
このとき根は除法であり、左の子は第一の基本ブロックである。
残りの列のすべての符号を反転させると、右の子が再帰的に復元される。
したがって、符号付き因子列、ゆえに有理関数から式が一意に定まる。
\end{proof}

\begin{theorem}\label{thm:main}
$F_1=F_2=1$ とする。このとき、
\[
 r_2=3,
 \qquad
 r_n=2F_{n+1}\quad(n\ge3)
\]
である。したがって、
\[
 r_2,r_3,r_4,r_5,r_6,r_7,r_8,\ldots
   =3,6,10,16,26,42,68,\ldots
\]
となる。
\end{theorem}

\begin{proof}
$r_2=3$ であることはすでに説明した。また、三変数の場合に表示した式の一覧から $r_3=6$ である。

葉が四つの場合、根が乗法である式は一つだけであり、左右に二葉の和を一つずつもつ。
根が除法である式の個数は、左のブロックが一つまたは二つの葉をもつ場合に応じて、
\[
 r_3+r_2=6+3=9
\]
である。したがって $r_4=10$ である。

次に $n\ge5$ とする。根は $+$ にも $\times$ にもなり得ない。
前者が許す葉は高々二つであり、後者が許す葉は高々四つだからである。したがって根は $/$ である。
左の子は葉 $x_1$ であって、$x_2,\ldots,x_n$ 上に任意の剛直な式を残すか、
または和 $x_1+x_2$ であって、$x_3,\ldots,x_n$ 上に任意の剛直な式を残す。
この二つの場合は互いに素であり、すべての場合を尽くす。よって、
\[
 r_n=r_{n-1}+r_{n-2}\qquad(n\ge5)
\]
である。$r_3=6=2F_4$ および $r_4=10=2F_5$ であるから、主張した公式が従う。
\end{proof}

直前の証明に現れた二つの選択肢は、フィボナッチ数の背後にあるおなじみのモノマー--ドミノ選択である。
左から読めば、少なくとも五つの変数をもつすべての剛直な式は、
一変数の分子または二変数の加法的な分子で始まり、その後に除法と、より小さな剛直な式が続く。
同値な言い方をすれば、$n\ge3$ に対して $r_n$ は、$n$ を $1$ と $2$ の和として表す組成の個数の二倍である。
ただし、構造的な漸化式そのものが始まるのは $n=5$ からであり、$n=3,4$ が初期値を与える。

\begin{corollary}\label{cor:gf}
非自明な剛直式の通常母関数は、
\[
 R(z)=\sum_{n\ge2}r_nz^n
     =\frac{3z^2+3z^3+z^4}{1-z-z^2}
\]
である。
\end{corollary}

\begin{proof}
$\Zcal$ を葉、$\Pcal$ を二つの葉を子にもつ加法節点とし、
$\Bcal=\Zcal\mathbin{\dot\cup}\Pcal$ とおく。
命題~\ref{prop:structure} から、次の曖昧さのない仕様が得られる。
\[
 \Rcal
 =\Pcal
 \mathbin{\dot\cup}
 \operatorname{Mul}(\Bcal,\Bcal)
 \mathbin{\dot\cup}
 \operatorname{Div}(\Bcal,\Zcal\mathbin{\dot\cup}\Rcal).
\]
葉を $z$ に置き換えると、
\[
 R=z^2+(z+z^2)^2+(z+z^2)(z+R)
\]
となり、これを $R$ について解けば、表示した有理関数を得る。
\end{proof}

\section{二つの精密化}

第一に、数え上げに各演算を記録させることができる。
加法、乗法、除法にそれぞれ重み $u,v,w$ を与え、対応する母関数を $R(z;u,v,w)$ とする。
\[
 B(z;u)=z+uz^2
\]
とおけば、重み付き構造仕様は
\[
 R=uz^2+vB(z;u)^2+wB(z;u)(z+R)
\]
となる。これを $R$ について解くと、
\[
 R(z;u,v,w)=
 \frac{(u+v+w)z^2+u(2v+w)z^3+u^2vz^4}
      {1-wz-uwz^2}
\]
を得る。したがって、通常のフィボナッチ型の分母は、重み $wz$ をもつ一葉の分子と、
重み $uwz^2$ をもつ二葉の加法的な分子との間で繰り返し選択することを記録している。

第二に、すべての式の中で剛直な式は稀である。
$n$ 個の葉をもつ平面完全二分木は $C_{n-1}$ 個あり、その内部頂点に演算を付す方法は $3^{n-1}$ 通りある。
したがって $n\ge3$ に対して、剛直な式の割合は
\[
 \frac{2F_{n+1}}{3^{n-1}C_{n-1}}
\]
であり、これは指数的に $0$ へ近づく。
$\varphi=(1+\sqrt5)/2$ とおくと、カタラン数とフィボナッチ数の標準的な漸近公式から、
\[
 \frac{r_n}{3^{n-1}C_{n-1}}
 \sim
 \frac{2\varphi^2\sqrt{\pi}}{\sqrt5}
 (n-1)^{3/2}\left(\frac{\varphi}{12}\right)^{n-1}
\]
を得る。

\section*{位置づけ}

冗長な括弧に関するいくつかの研究では、優先順位と結合規則のもとで構文を保存する文法や、
構文木から中置式を出力するアンパーシング・アルゴリズムが扱われている
\cite{Ramsey,Bates,Lotfallah}。本稿における比較は意味論的である。
すなわち、二つの表示が同じ有理関数を表すとき、それらを同一視する。
算術式に関する別の数え上げ問題では、あらかじめ指定された括弧の必要性ではなく、
有理関数の同値類が扱われている \cite{StosicDamnjanovicRandelovic,Cohen}。
本稿の個数が小さいのは、すべての括弧が意味をもつことを要求すると、右方向に延びる剛直な商構造が強制されるからである。
フィボナッチ漸化式は、すでに根における選択に現れている。

OEIS の記法では、
\[
 r_n=A118658(n+2)\qquad(n\ge3)
\]
である。実際、A118658 は $m>0$ に対して $a(m)=2F_{m-1}$ を満たす \cite{OEIS}。
定理~\ref{thm:main} は、この数列に括弧による解釈を与える。
OEIS A155069 は別のモデルを記録している。それは $+$ と $-$ を用いた、括弧が最小限に付された中置文字列を数える。
一方、本稿では、あらかじめ指定された各内側の括弧が必要となる完全括弧付き木を数える。
A155069 のシュレーダー型の増大は、本稿において乗法と右側の除法によって強制されるフィボナッチ構造と対照的である。

\begin{thebibliography}{9}
\footnotesize
\setlength{\itemsep}{0pt}

\bibitem{Bates}
R. M. Bates,
Distinguishing redundant parentheses by syntax,
\emph{Internat. J. Comput. Math.} \textbf{82} (2005), 951--967.

\bibitem{Cohen}
B. Cohen,
The number of non-isomorphic arithmetic expressions that can be constructed
using $+,-,\times$, and $/$,
arXiv:2602.18522 (2026).

\bibitem{Lotfallah}
W. B. Lotfallah,
Characterizing unambiguous precedence systems in expressions without
superfluous parentheses,
\emph{Internat. J. Comput. Math.} \textbf{86} (2009), 1--20.

\bibitem{Ramsey}
N. Ramsey,
Unparsing expressions with prefix and postfix operators,
\emph{Software: Practice and Experience} \textbf{28} (1998), 1327--1356.

\bibitem{StosicDamnjanovicRandelovic}
I. Sto\v{s}i\'c, I. Damnjanovi\'c, and \v{Z}. Ran{\dj}elovi\'c,
Counting the number of inequivalent arithmetic expressions on $n$ variables,
\emph{Filomat} \textbf{39} (2025), 949--962.

\bibitem{OEIS}
N. J. A. Sloane and The OEIS Foundation Inc.,
The On-Line Encyclopedia of Integer Sequences, A118658 and A155069.

\end{thebibliography}

\end{document}
